\par\Needspace{6\baselineskip}\medskip
\hypertarget{native-111-article082}{}\pdfbookmark[2]{82.}{article-082}
\begin{center}82.\end{center}
Supponamus denique, tres radios vectores $r$, $r'$, $r''$ datos esse, qui cum recta ad lubitum e Sole in plano orbitae ducta faciant angulos $N$, $N'$, $N''$. Habebuntur itaque, retentis signis reliquis, aequationes:
\[
\begin{aligned}
&\frac{p}{r}=1+e\cos(N-\Pi)\\[4pt]
\mathrm{I.}\qquad&\frac{p}{r'}=1+e\cos(N'-\Pi)\\[4pt]
&\frac{p}{r''}=1+e\cos(N''-\Pi),
\end{aligned}
\]
e quibus $p$, $\Pi$, $e$ pluribus modis diversis elici possunt. Si quantitatem $p$ ante reliquas computare placet, multiplicentur tres aequationes I resp. per $\sin(N''-N')$, $-\sin(N''-N)$, $\sin(N'-N)$, fietque additis productis per lemma I art. 78
\[
p=\frac{\sin(N''-N')-\sin(N''-N)+\sin(N'-N)}{\dfrac1r\sin(N''-N')-\dfrac1{r'}\sin(N''-N)+\dfrac1{r''}\sin(N'-N)}\mathpunct{\cdot}
\]
Haec expressio propius considerari meretur. Numerator manifesto fit
\[
\begin{aligned}
&=2\sin\tfrac12(N''-N')\cos\tfrac12(N''-N')\\
&\qquad-2\sin\tfrac12(N''-N')\cos(\tfrac12N''+\tfrac12N'-N)\\
&=4\sin\tfrac12(N''-N')\sin\tfrac12(N''-N)\sin\tfrac12(N'-N).
\end{aligned}
\]
Statuendo porro
\[
\begin{aligned}
r'r''\sin(N''-N')&=n\\
rr''\sin(N''-N)&=n'\\
rr'\sin(N'-N)&=n'',
\end{aligned}
\]
patet $\tfrac12n$, $\tfrac12n'$, $\tfrac12n''$ esse areas triangulorum inter radium vectorem secundum et tertium, inter primum et tertium, inter primum et secundum. Hinc facile perspicietur, in formula nova
\SourceLeaf{112}
\[
p=\frac{4\sin\tfrac12(N''-N')\sin\tfrac12(N''-N)\sin\tfrac12(N'-N)\mathbin{.}rr'r''}{n-n'+n''}
\]
denominatorem esse duplum areae trianguli inter trium radiorum vectorum extremitates i. e. inter tria corporis coelestis loca in spatio contenti. Quoties haec loca parum ab invicem remota sunt, area ista semper erit quantitas perparva et quidem ordinis tertii, siquidem $N'-N$, $N''-N'$ ut quantitates parvae ordinis primi spectantur. Hinc simul concluditur, si quantitatum $r$, $r'$, $r''$, $N$, $N'$, $N''$ una vel plures erroribus utut levibus affectae sint, in determinatione ipsius $p$ errorem permagnum illinc nasci posse; quamobrem haecce ratio orbitae dimensiones eruendi magnam praecisionem nunquam admittet, nisi tria loca heliocentrica intervallis considerabilibus ab invicem distent.

Ceterum simulac semiparameter $p$ inventus est, $e$ et $\Pi$ determinabuntur e combinatione duarum quarumcunque aequationum I per methodum art. 79.
